\section{Evidence Contracts}
\label{sec:contracts}

\subsection{Plans and execution records}
\label{sec:algebra}

Let $\mathcal{S}$ contain relational tables, text collections with sparse
and dense indexes, vector stores, and structured APIs. A physical plan
$P$ is a DAG over the typed operators
\[
\textsf{Retrieve}_{s,k}\mid\textsf{Fuse}\mid\textsf{Rerank}_m\mid
\textsf{Filter}_\phi\mid\textsf{Generate}_g\mid\textsf{Repair}.
\]
\textsf{Retrieve} implements SQL selection, BM25, dense ANN, or an API
call; \textsf{Fuse} combines rankings; \textsf{Rerank} applies model $m$;
\textsf{Filter} enforces an evidence predicate; \textsf{Generate} invokes
generator $g$; and \textsf{Repair} revises an answer after verification.
Executing $P$ on query $q$ returns an answer, an evidence set with citation
links, and an execution record containing cost, latency, evidence age, and
operator-level statistics.

Two such plans need not return the same answer. We therefore optimize over
a \emph{policy} $\pi$, which maps a query and intermediate execution records
to a schedule of one or more plans. A static plan, a direct router, and a
progressive ladder are all policies.

\subsection{Contract and optimization problem}

\begin{definition}[Evidence contract]
\label{def:contract}
For a query distribution $\mathcal{D}$, an evidence contract is
$C=\langle(\ell_j,\alpha_j)_{j=1}^{k},\delta\rangle$. Each bounded loss
$\ell_j(\pi,q)\in[0,1]$ records a violation, $\alpha_j$ is its admissible
population rate, and $\delta$ is the probability budget for an erroneous
deployment decision.
\end{definition}

Examples include
\[
\begin{split}
\ell_{\mathrm{qual}}&=\mathbf{1}[\mathrm{quality}<\tau_q],\qquad
\ell_{\mathrm{cite}}=\mathbf{1}[\mathrm{citationRecall}<\tau_c],\\
\ell_{\mathrm{fresh}}&=\mathbf{1}[\max_{e\in E(q)}\mathrm{age}(e)>\Delta_f],
\quad
\ell_{\mathrm{lat}}=\mathbf{1}[\mathrm{latency}>B_\ell].
\end{split}
\]
The risks and mean cost are
$R_j(\pi)=\E_{q\sim\mathcal D}[\ell_j(\pi,q)]$ and
$c(\pi)=\E_{q\sim\mathcal D}[c(\pi,q)]$. The optimizer solves
\begin{equation}
\label{eq:problem}
\min_{\pi\in\Pi} c(\pi)
\quad\text{s.t.}\quad R_j(\pi)\leq\alpha_j,\quad j=1,\ldots,k .
\end{equation}
Because $\mathcal D$ is unknown, the system receives a finite calibration
sample. Its output is itself random.

\begin{definition}[Certified deployment]
\label{def:certified}
A selection procedure returning $\hat\pi$ from calibration data is
$(C,\delta)$-certified if
\[
\Prob\!\left(R_j(\hat\pi)\leq\alpha_j\ \text{for all }j\right)
\geq 1-\delta ,
\]
where probability is over the calibration sample.
\end{definition}

The contract governs the audited service-level quantity, not individual
answers. Its loss may use gold labels, a benchmark evaluator, evidence
timestamps, measured latency, or human audits. The optimizer does not
assume that a learned score is accurate; it requires only bounded
calibration losses. This separation is important: learned estimates can
improve the cost of the selected policy, but only direct testing may
declare the policy feasible.
